% An independent 100-operation encoding, with all Pell equations retained.
\section{Halving the coefficient and sharing the geometric factor}
\label{sec:half-center}

The fixed-four construction admits another saving, independent of the
Pell change in the preceding section. Starting from
Theorem~\ref{thm:best101}, replace its positive coefficient and third
block by
\begin{equation}\label{eq:half-center-block}
 \Omega=2\lambda-e>0,\qquad
 S_3=\theta\lambda q-\Omega\mathcal C^2=\sigma>0,
 \qquad \mathcal C=x+g,\quad\theta=B-4.
\end{equation}
The factor $\theta\lambda$ can now be shared with the geometric
equation. To accommodate the smaller coefficient, rebuild the fixed
index by pairing ordinary equation tests and placing the special
target first. No packed bound or Pell equation is removed.

\subsection{Paired rows and a first special target}

Use the finite primitive circuit construction of
Section~\ref{sec:fixed-four}, with distinct copied operands and fresh
outputs, all separate from the homogenizing coordinate $\delta$.
For every multiplication, addition, copy, final equality, or zero
residual
\[
 XY-W\delta,\qquad (X+Y-W)\delta,\qquad
 (X-Y)\delta,\qquad X\delta,
\]
include both $F$ and $-F$ among the ordinary rows. Omit trivial
identical equalities. Replace each unit gate $X=1$ by two unpaired
rows
\begin{equation}\label{eq:half-center-unit}
 F_{X,1}=X\delta-\delta^2,\qquad
 F_{X,2}=X\delta-X^2.
\end{equation}
Include the unpaired guard $F_{\mathrm g}=u\delta-x^2$, with $u$
fresh. Every ordinary row has target
\begin{equation}\label{eq:half-center-target}
 G_j=2F_j+\delta^2.
\end{equation}
The special target $G_0=\delta^2$ is placed \emph{first}, before all
ordinary targets.

After division by the monomial coefficient in $\mathcal C(T)^2$,
the coefficient alphabet is still $\{-2,-1,0,1\}$. Paired primitive
rows have cross-term quotients $\pm1$ and the added square
$\delta^2$ has quotient $1$. The two unit targets are
\[
 2X\delta-\delta^2,\qquad
 2X\delta-2X^2+\delta^2,
\]
and the guard has the latter coefficient pattern. The special
target has its single coefficient $1$. Thus the pairing does not
enlarge the digit alphabet. At a genuine circuit solution with
$\delta=1$, every ordinary $F_j$ is zero, including both unit rows;
choose $u=x^2$ for the guard.

\subsection{The rebuilt fixed index and the absence of an initial carry}

Apply the dynamic support construction of Section~\ref{sec:fixed-four}
to this finite list. With $m$ positive-weight coordinates and $s$
targets, write
\begin{align*}
 v_i&=3^i\quad(0\le i<m),& M&=3^{m-1},& d_0&=4M+1,\\
 t_0&=(s+1)d_0+2M,& t_j&=t_0+jd_0,& t_*&=2sd_0+2M,\\
 K&=t_*+1,& L&>3K+2,
\end{align*}
where $L$ is a power of two, $x$ has weight zero, and $\delta$ has
weight $v_0=1$. The special target is at $t_0$. Include the usual
dummy coordinate at each $t_j$. Monomial uniqueness, separation of
the intervals $[t_j-2M,t_j]$, and $2t_0-2M>t_*$ give
\begin{equation}\label{eq:half-center-coefficient}
 [T^{t_j}]\,\mathcal D(T)\mathcal C(T)^2=G_j.
\end{equation}
In particular, no term involving a dummy contributes to any target.

Define the fixed polynomials and parameters by
\begin{align*}
 \ell_0(T)&=\sum_{i=0}^{m-1}T^{v_i}+\sum_{j=0}^{s-1}T^{t_j},\\
 e_0(T)&=2\sum_{j=0}^{K-1}T^j+\mathcal D(T),&
 V&=\ell_0(4)+e_0(4)4^L.
\end{align*}
The coefficients of $e_0$ are in $\{0,1,2,3\}$, with unit digit
$2$. Put $c_*=m+s+1$, counting all code coordinates including the
input, and choose a power of two
\begin{equation}\label{eq:half-center-index-bound}
 H>\max\{2\,4^{2L+1},\ 4^{t_*+3}\,4c_*^2,\ 3L,\ 16\}.
\end{equation}
The resulting $(V,H,L)$ is an admissible index. Its construction is
independent of the queried input and introduces no additional
unknown in the final system.

The first target has a stronger support property than the others.
Its sole coefficient in $\mathcal D$ is $+1$ at $t_0-2$.
Every later row has its entire coefficient support above $t_0$,
because
\[
 t_1-2M=t_0+d_0-2M>t_0.
\]
Consequently every raw coefficient of
$\mathcal D(T)\mathcal C(T)^2$ through position $t_0$ is
nonnegative. After coordinate decoding these coefficients are below
$B$, as proved below. Hence the incoming carry at the special target
is exactly zero. This argument assumes neither $\delta=1$ nor any
ordinary equation.

\subsection{Bounds before the Pell argument}

Retain $B=Hb^2$, $b=x+\beta$, $n=q^8$, and
\begin{equation}\label{eq:half-center-coding}
 \lambda(B-1)=q^2-1,\quad \theta+4=B,\quad
 S_2=\ell+eq=V+t\theta,\quad S_2+\alpha=q^2.
\end{equation}
All supplied unknowns, including $\Omega,\sigma,\alpha$, remain
positive integers. Keep the packing
\begin{align}
 S&=g+q^2(S_2+q^2S_3),\notag\\
 T+1&=q^2(1+\theta\lambda)-(b-1)\ell+\theta\ell q^4,\notag\\
 r&=S(n^2-n)+(T+1)(n^2-1).\label{eq:half-center-packing}
\end{align}

The packed bound gives $e<q$ and $\ell<q^2$. Geometry and
\eqref{eq:half-center-index-bound} give
\[
 B\le q^2,\qquad q>4b\ge8,\qquad 3L\le B\le n.
\]
Since $\Omega\ge1$ and $S_3>0$, the new offset gives
\begin{equation}\label{eq:half-center-preliminary-C}
 \mathcal C^2<\theta\lambda q<B\lambda q<3q^3<q^4.
\end{equation}
Thus $g<\mathcal C<q^2$ and $0<S_3<q^4$. The mask $T+1$ is
unchanged; its first block borrows fewer than $b$ units from the
middle block, which remains positive since $\theta=B-4>b$.
The same width-$(2,2,4)$ bounds therefore give
\begin{equation}\label{eq:half-center-packing-bounds}
 0<S<n,\qquad 0<T+1\le n,\qquad
 n\le r\le2n^3-2n^2<2n^3.
\end{equation}

These bounds have been established before any Pell equation is
used. In particular $n,r\ge64$, $U=w n^2\ge n^2$ and
$Y_P=s n^2\ge n^2$. The main parameter satisfies
$a=Y_P(U+1)>n^4>2r+1$, and the first index modulus satisfies
$UY_P\ge n^4>r+1$. Moreover $L\ge2$ and
$3L\le B\le n\le r$. Thus every preliminary hypothesis of the
common-witness Pell proof in Section~\ref{sec:common-witness} is
available. All its equations are retained. That proof yields powers
of two $b,B,q$, the exact equality $q=B^L$, and
\begin{equation}\label{eq:half-center-binomial}
 n^2\mid\binom{2r}{r}.
\end{equation}
No internal target equation has entered this argument.

Now $\lambda>q$ and $e<q$, so $\Omega=2\lambda-e>\lambda$.
Positivity in \eqref{eq:half-center-block} improves the code bound to
\begin{equation}\label{eq:half-center-post-C}
 \mathcal C^2<\theta q<Bq<q^2,
 \qquad g<\mathcal C<q.
\end{equation}
The original factor one-half in the bound for $\mathcal C^2$ is no
longer asserted; the strict inequality $\mathcal C<q$ is what the
subsequent decoding needs.

\subsection{Canonical decoding and the revised high window}

By \eqref{eq:half-center-packing-bounds} and
\eqref{eq:half-center-binomial}, the packed binary no-carry tests
hold. Let $d=\lfloor(b-1)\ell/q^2\rfloor<b$. Subtracting $d$
from $(B-4)\lambda$ changes only its unit digit. The middle mask
therefore bounds all nonunit digits of $S_2$ by three and its unit
digit by $3+d$. It follows that
\[
 \ell\le S_2<\frac{4q^2}{B-1}+b.
\]
Together with $B=Hb^2$ and $H>16$, this forces
$(b-1)\ell<q^2$, hence $d=0$. The base-four evaluation lemma now
applies to the fixed polynomial defining $V$: its degree is below
$2L$, all digits are below four, and the first term of
\eqref{eq:half-center-index-bound} supplies the evaluation bound.
Consequently
\begin{equation}\label{eq:half-center-alias}
 \begin{gathered}
 S_2=\ell_0(B)+e_0(B)q,\\
 \ell=\ell_0(B)+a_{\mathrm{al}}q,\qquad
 e=e_0(B)-a_{\mathrm{al}},\qquad
 0\le a_{\mathrm{al}}<e_0(B).
 \end{gathered}
\end{equation}

The first mask and $g<q$ decode the low positions specified by
$\ell_0$. The unit digit of $g$ is zero; every coordinate is a
nonnegative integer below $b$, and the unit digit of $\mathcal C$
is $x$. Dummy digits occur only at row positions. Write
\[
 A_{\mathcal C}=\mathcal C(1)^2,
 \qquad J_{\mathcal C}=2A_{\mathcal C}.
\]
Then $\mathcal C(1)<c_*b$, and
\eqref{eq:half-center-index-bound} gives $J_{\mathcal C}<B/4$.
Also $e,\mathcal C<B^K$ and $L>3K+2$ give $e\mathcal C^2<q$.
Since $g>0$ and $x>0$, the decoded coefficient sum is at least two,
so
\begin{equation}\label{eq:half-center-coefficient-sum}
 A_{\mathcal C}\ge4.
\end{equation}
This holds before identifying $\delta$, even if the nonzero digit
of $g$ occupies a dummy position.

Write
\[
 2\lambda\mathcal C^2=qV_1+V_0,\qquad0\le V_0<q.
\]
Its polynomial coefficients are at most $J_{\mathcal C}<B$, so
they are its base-$B$ digits. The stable convolution interval
$L,\ldots,L+K$ has all digits equal to $J_{\mathcal C}$, because
$L>3K+2$. Exact division of \eqref{eq:half-center-block} by $q$
therefore gives
\[
 \left\lfloor\frac{S_3}{q}\right\rfloor
 =\theta\lambda-V_1+\epsilon,\qquad
 \epsilon=\left\lfloor
 \frac{e\mathcal C^2-V_0}{q}\right\rfloor\in\{-1,0\}.
\]
Every digit of $\theta\lambda$ through position $2L-1$ is $B-4$.
Subtracting $V_1$ creates no borrow in the window: its first digit
$B-4-J_{\mathcal C}+\epsilon$ is positive, as are all subsequent
digits. Thus
\begin{equation}\label{eq:half-center-high-digits}
 (S_3)_L=B-4-J_{\mathcal C}+\epsilon,\qquad
 (S_3)_{L+j}=B-4-J_{\mathcal C}\quad(1\le j\le K).
\end{equation}

The high part of the third mask is
$X=(B-4)a_{\mathrm{al}}<B^{K+1}$. Binary nonoverlap bounds each
digit of $X$ by
\[
 J_{\mathcal C}+4=2A_{\mathcal C}+4\le4A_{\mathcal C},
\]
using \eqref{eq:half-center-coefficient-sum}. If $F_X$ is its digit
polynomial, then $B-4\mid F_X(4)$, whereas
\[
 0\le F_X(4)\le
 4A_{\mathcal C}\frac{4^{K+1}-1}{3}<\frac B{12}<B-4.
\]
The strict bound follows from
\eqref{eq:half-center-index-bound}, without enlarging $H$.
Therefore $F_X(4)=0$, $X=0$, and $a_{\mathrm{al}}=0$.
Both supplied codes are canonical.

\subsection{The low target tests and positive necessity}

With the canonical codes, the raw coefficients of $S_3$ below $K$
are those of $\mathcal D(T)\mathcal C(T)^2$. The positive offset
starts at $L>K$, and the missing baseline tail starts at $K$.
Every raw coefficient has absolute value at most
$2A_{\mathcal C}<B/4$, so every incoming carry is zero or minus one.
At an ordinary target put $v=2F+\delta^2$. The mask $B-4$ permits
exactly the digits $0,1,2,3$ and is equivalent to
\begin{equation}\label{eq:half-center-allowed}
 0\le v+\epsilon\le3,\qquad\epsilon\in\{-1,0\}.
\end{equation}
A negative value would normalize to a digit greater than three.

At the first special target there is no incoming carry, by its
support property. The target is $\delta^2$, so
$0\le\delta^2\le3$, forcing $\delta\in\{0,1\}$. If $\delta=0$,
the guard has $v=-2x^2<0$ and cannot pass. Thus $\delta=1$.
For either possible carry, \eqref{eq:half-center-allowed} now says
\[
 0\le2F+1+\epsilon\le3
 \quad\Longleftrightarrow\quad F\in\{0,1\}.
\]
Paired rows $F$ and $-F$ force $F=0$. The unit rows give
\[
 X-1\in\{0,1\},\qquad X-X^2\in\{0,1\}.
\]
The first restricts $X$ to $\{1,2\}$, and the second restricts it
to $\{0,1\}$, so $X=1$. The guard need not vanish in sufficiency:
it excludes $\delta=0$, and its coordinate $u$ is absent from the
original circuit. Every original gate and equality has now been
recovered exactly at the input $x$. The low coordinate tests are
automatic because $\mathcal D$ starts above their positions.

Conversely, take a solution of the represented system, extend it by
its circuit values, and choose $\delta=1$, $u=x^2$, and zero dummy
coordinates. Choose a power-of-two $b$ larger than every coordinate
and use the canonical codes and $q=B^L$. Every ordinary $F$ is
zero. The special target has raw coefficient one and carry zero;
each ordinary target has raw coefficient one and actual digit zero
or one. These digits are allowed, and the low variable tests vanish.
The first two masks hold as in Section~\ref{sec:fixed-four}.

All positive slack witnesses exist. In particular, the unit digit
of $e_0$ is two, the positive $\delta$ digit makes $g>0$, and
$S_2<q^2$. The packed-congruence quotient is positive because
$P(B)>P(4)$ for the nonconstant polynomial
$P(T)=\ell_0(T)+e_0(T)T^L$ with nonnegative coefficients.
Also $\Omega=2\lambda-e>\lambda>0$. Since
$\mathcal C<B^K$ and $L>3K+2$,
\[
 2\mathcal C^2<q\le\theta q,
 \qquad
 S_3=\lambda(\theta q-2\mathcal C^2)+e\mathcal C^2>0.
\]
The three masks give \eqref{eq:half-center-binomial}, and their
packing satisfies \eqref{eq:half-center-packing-bounds}. The complete
positive Pell witness construction of Section~\ref{sec:common-witness}
therefore applies to the newly calculated $r$. This proves necessity
for the rebuilt fixed index.

\subsection{Seven shared instructions in place of eight}

The predecessor uses eight instructions for the relevant quantities:
\[
 \begin{gathered}
 L_2=\lambda+q^2,\quad \Lambda=B\lambda,\quad R_2=\Lambda+1,\\
 z_\lambda=4\lambda,\quad e_2=2e,\quad
 t_2=z_\lambda-e_2,\\
 T_{\mathrm{coef}}=L_2-z_\lambda,\qquad P_2=\Lambda q.
 \end{gathered}
\]
Use instead the seven instructions
\begin{equation}\label{eq:half-center-seven}
 \begin{gathered}
 U_\lambda=\theta\lambda,\qquad \lambda_2=2\lambda,\qquad
 T_{\mathrm{coef}}=U_\lambda+1,\\
 \lambda_3=\lambda_2+\lambda,\qquad
 R_2=T_{\mathrm{coef}}+\lambda_3,\\
 t_2=\lambda_2-e,\qquad P_2=U_\lambda q.
 \end{gathered}
\end{equation}
Test $q^2=R_2$ and $t_2=\Omega$. In conjunction with
$\theta+4=B$, the first test is exactly
$\lambda(B-1)=q^2-1$. All other arithmetic instructions remain,
including the packed upper bound and $S_3=P_2-t_2\mathcal C^2$.

\begin{theorem}\label{thm:best100-encoding}
Every recursively enumerable subset of the positive integers has an
admissible fixed index $(V,H,L)$ for a system with 34 positive unknowns
and 22 equations admitting a straight-line certificate of 100
operations: 55 multiplications and 45 additions. The supplied
parameters are $(x,V,H,L)$; fixed numerals are free.
\end{theorem}
\begin{proof}
The construction above proves both positive-domain implications.
The replacement \eqref{eq:half-center-seven} has three multiplications
and four additions, instead of four multiplications and four
additions, and every other instruction is retained. Hence the
101-operation certificate loses one multiplication.

The exact verification is
\texttt{round25\_1980\_half\_center\_certificate.py} and its JSON
receipt. It checks the full 100-instruction primitive list and all
22 polynomial residuals. If $F_3$ is the original geometric residual
and $F_4=\theta+4-B$, the new geometric test has residual
$F_3-\lambda F_4$. The new packing and positive-$S_3$ source equations
use $\theta$ explicitly; they need no geometric correction.
The existing signed Pell and second-exponent corrections are also
verified. Thus the count concerns an explicit certificate for the
proved system; no optimality is claimed.
\end{proof}
